\subsection{Experimental inputs and current status}

    The next scientific step is validation on independently collected
    experimental units. This step has been prepared, not performed: the
    checkout contains only the development CSV already used in the preceding
    investigations. Its time and response units, experimental unit definition,
    and common experimental conditions remain unconfirmed. Three trial
    columns cannot by themselves establish three independent experiments.
    The metadata record retains these unknowns explicitly; inherited figure
    labels are not used to fill them.

    Table~\ref{tab:prospective-readiness} summarizes the current inputs, and
    Table~\ref{tab:prospective-design} records the draft comparison. The
    proposed origins and horizon retain original numerical time units and
    require confirmation against the experimental schedule. The four model
    candidates, priors, half/double noise sensitivities, and MCMC gates are
    retained from the structural comparison. No candidate is selected by
    the new outcomes, which have not been supplied. The draft is not a
    completed prospective registration or a new calibration result.

% Generated by scripts/build_prospective_tables.py; do not edit manually.
\begin{table}[htbp]
\centering
\caption{Current validation readiness. Declared new units are not inferred from trial columns or file names.}
\label{tab:prospective-readiness}
\begingroup\small\setlength{\tabcolsep}{4pt}
\begin{tabular}{ll}
\toprule
Item & Value \\
\midrule
Development time points & 15 \\
Development trial series & 3 \\
Development scalar responses & 45 \\
Declared new independent units & 0 \\
Unresolved experimental declarations & 9 \\
Protocol ready to freeze & No \\
\bottomrule
\end{tabular}
\endgroup
\end{table}
\begin{table}[htbp]
\centering
\caption{Draft prospective comparison, pending experimental metadata and schedule confirmation. Times retain original numerical units.}
\label{tab:prospective-design}
\begingroup\small\setlength{\tabcolsep}{4pt}
\begin{tabular}{ll}
\toprule
Item & Value \\
\midrule
Candidate models & 4 \\
Forecast origins & 108, 120, 132, 144, 156 \\
Forecast horizon & 12 \\
Observation interval level & 0.90 \\
Primary score & Central interval score \\
Point forecast & Posterior predictive median \\
Aggregation & Equal weights for independent units \\
\bottomrule
\end{tabular}
\endgroup
\end{table}


\subsection{Separate forecasts from outcome access}

    Before target outcomes are inspected, specify the independent units,
    collection provenance, measurement units, forecast origins, horizons,
    candidate models, priors, diagnostics, and scoring rules. A local
    protocol snapshot copies the specification, experimental metadata, and
    relevant inference sources, recording hashes and a UTC timestamp. It
    refuses to overwrite an existing snapshot. This supports change
    detection; it does not establish external preregistration or prove the
    time of first outcome access. The snapshot operation does not inspect
    the new prefix observations, which may not yet exist.

    At each origin, fit each candidate to that unit's available prefix and
    archive its predictive distribution before inspecting the corresponding
    target responses. This is a within-trajectory prediction study on new
    units. It does not transfer a single historical fitted capacity to all
    experiments or introduce hierarchical pooling. Prefix checks reject
    future rows, altered earlier observations, and exact copies of
    development prefixes relabeled as new data. They cannot establish
    experimental independence without the experimental record.

    An outcome-free forecast interface constructs an equally weighted
    posterior mixture of joint Gaussian log responses. Component covariance
    matrices retain the modeled dependence across times and trials, including
    the difference between shared and independent discrepancies. A separate
    scoring interface receives the new observed responses and computes
    central 90 percent observation intervals, interval scores, marginal
    predictive densities, and the joint mixture density. Density scores
    retain the response transformation Jacobian. Point errors use the
    predictive median; parameter integration is retained in the mixture CDF
    and densities. Gaussian density
    oracles verify these calculations. The distinction between marginal
    and joint scores is retained because their earlier model rankings
    disagreed. Interval and density scores assess predictive quality beyond
    coverage alone \citep{gneiting2007proper}.

\subsection{Predictive moments and the point forecast}
\label{sec:predictive-moments}

    Before freezing the draft, a tail check identifies a limitation of the
    response-scale means used in the preceding finite-subset calculations.
    Let \(z_b=\log b\) have its declared normal prior and let \(N\) be the
    number of training scalar responses. Keeping other parameters in a
    compact interior set, as \(z_b\to\infty\),
    \[
        \log L(z_b)=-N z_b+O(1),\qquad
        \log\mathbb E[Y_*\mid z_b,\theta,D]
        =\tfrac12 e^{2z_b}+O(1).
    \]
    The first identity follows from the observation covariance
    \(b^2I+O(1)\); the second follows from the conditional lognormal
    distribution, whose measurement variance is \(b^2\). Thus the
    unnormalized first-moment integrand has log asymptotic
    \[
        \tfrac12 e^{2z_b}
        -\frac{(z_b-\mu_b)^2}{2d_b^2}
        -N z_b+O(1)\longrightarrow+\infty.
    \]
    The bound holds on a positive-volume compact set of the other parameters,
    so integration over their joint posterior also gives
    \(\mathbb E[Y_*\mid D]=\infty\). Increasing the number of observations
    only changes the linear term. The likelihood tail cannot overcome this
    double-exponential term under the declared lognormal prior for \(b\).
    This applies to both sharing assumptions and both temporal kernels.
    The posterior itself is proper: the observation likelihood is bounded
    above by a data-dependent constant times \(b^{-N}\), whose prior
    expectation is finite for lognormal \(b\). Its positive normalizing
    constant is therefore finite. Existence of this distribution and
    existence of its response moments are separate requirements.

    Finite posterior-subset means are finite numerical forecast summaries;
    their apparent stability does not establish a finite full-posterior
    response mean. The earlier mean-based RMSE and bias numbers are retained
    with that qualification. Central predictive quantiles and predictive
    densities remain well-defined functionals of the proper distribution.
    The prospective API consequently uses the posterior predictive median
    for point errors, not the population logistic median \(m_\theta\).
    Its RMSE is a descriptive error of this specified median forecast, not
    a claim of finite posterior quadratic risk. A physically justified prior
    redesign would be necessary if finite response means were required.
    No upper noise bound is invented while physical units are unknown.
    Section~\ref{sec:tempered-noise} investigates a separate smooth
    prior modification; the prospective draft itself is not changed.

    A reproducible log-scale tail check in
    \path{results/prospective/predictive_moment_tail.csv} verifies the
    \(-N z_b\) likelihood asymptotic numerically. It supports the independent
    analytic argument; a finite grid alone cannot prove divergence.

\subsection{Independent units determine comparison weights}

    Let \(G\) be the number of independently collected units and
    \(\mathcal I_{g,h}\) the target responses for unit \(g\) at horizon \(h\).
    For a pointwise loss \(L\), the proposed summary is
    \[
        \bar L_{m,h}
        =\frac{1}{G}\sum_{g=1}^{G}
            \frac{1}{|\mathcal I_{g,h}|}
            \sum_{i\in\mathcal I_{g,h}}L_{m,g,h,i}.
    \]
    Thus a unit with more repeated measurements does not receive greater
    weight in the across-unit comparison. Coverage uses indicator losses;
    RMSE is the square root of the resulting mean squared-error summary,
    rather than an average of unit RMSEs. Model comparisons require identical
    unit and target identities at each horizon. Joint NLL is retained per
    forecast block and summarized within and across units separately; it is
    not reconstructed from marginal scores.

    Report each unit and horizon before aggregating. Paired model differences
    retain the same unit and forecast targets. Any future resampling-based
    uncertainty assessment must preserve entire independent units and model
    pairings. Treating longitudinal response rows as independent experiments
    would change the sampling unit. The present utilities do not generate
    calibration confidence intervals or assert statistical significance.

    The readiness checker, protocol snapshot, forecast-moment, and scoring
    APIs are implemented and tested. A complete automated new-data fitting
    and forecast-archiving runner is not yet supplied. Synthetic fixtures
    establish computational identities, not new kefir validation. Confirmed
    experimental metadata, additional independent trajectories, and a
    documented new-data run remain necessary. REV-07 and REV-10 therefore
    remain scientifically open.
